\documentstyle[psfig,proc,sectspaceeq]{article}
\pagestyle{empty}
%\input define
\textheight=24cm
\textwidth=17.7cm
\oddsidemargin=-.4cm
\columnsep=1.5cm
\topmargin=-2.716cm
%\def\thesection{\arabic{section}.}
%\def\thesubsection{\thesection\arabic{subsection}.}
%\def\theequation{\arabic{equation}}
%\def\figurename{\small Figure}
%******************************************************
%******From numdiff_SEPT8.tex

\newcommand{\r}{\hbox{\sl l\kern-.2em R}}
\def\thesection{\arabic{section}.}
\def\theequation{\thesection\arabic{equation}}
\def\theenumi{\roman{enumi}}\def\labelenumi{(\theenumi).}
\def\theenumii{\alph{enumii}}\def\labelenumii{(\theenumii)}
\newtheorem{theorem}{Theorem}[section]
\newtheorem{proof}{Proof:}
\def\theproof{}
\def\thetheorem{\thesection\arabic{theorem}}

\def\setR{\hbox{$l\! \! R$}}
\def\setN{\hbox{$l\! \! N$}}
\def\CN{{\cal N}}

\def \BA {{\bf A}}
\def \BB {{\bf B}}
\def \BC {{\bf C}}
\def \BE {{\bf E}}
\def \BF {{\bf F}}
\def \BG {{\bf G}}
\def \BH {{\bf H}}
\def \BI {{\bf I}}
\def \BJ {{\bf J}}
\def \Bk {{\bf k}}
\def \BK {{\bf K}}
\def \BL {{\bf L}}
\def \BM {{\bf M}}
\def \BN {{\bf N}}
\def \BQ {{\bf Q}}
\def \BP {{\bf P}}
\def \BQ {{\bf Q}}
\def \BR {{\bf R}}
\def \BS {{\bf S}}
\def \BV {{\bf V}}

%\def \CN   {{\eusm N}}

\def\thesubsubsection{\thesubsection\arabic{subsubsection}.}

%******************************************************

\makeatletter
\def\section{\@startsection{section}{1}{\z@}{\sectionspaceb
plus .5em minus .1em}{\sectionspacea plus 0em minus
.1em}{\centering\normalsize\bf}}

\def\subsection{\@startsection{subsection}{2}{\z@}{\sectionspaceb plus .5em minus
.1em}{\sectionspacea plus 0em minus .1em}{\normalsize\bf}}

\def\subsubsection{\@startsection{subsubsection}{3}{\z@}{\sectionspaceb plus .5em
minus .1em}{\qquad}{\underline{\normalsize\bf}}}

\def\caption{\refstepcounter\@captype \@dblarg{\@caption\@captype}}

%\long\def\@makecaption#1#2{\vskip 10pt \setbox\@tempboxa\hbox{#1: #2}
%\ifdim \wd\@tempboxa >\hsize #1: #2\par \else \hbox
%to \hsize{\hfil\box\@tempboxa\hfil} \fi}

\long\def\@makecaption#1#2{\hbox{#1: #2}}

\long\def\@caption#1[#2]#3{\addcontentsline{\csname
  ext@#1\endcsname}{#1}{\protect\numberline{\csname 
  the#1\endcsname}{\ignorespaces #2}}\par
  \begingroup
    \@parboxrestore
    \normalsize
    \@makecaption{\csname fnum@#1\endcsname}{\ignorespaces #3}\par
  \endgroup}

\makeatother

\begin{document}

\small

\baselineskip=12pt
\parskip=0in


\title{\vspace*{.5cm}
\large\bf Interpolation and Numerical Differentiation \\
for Observer Design}

\author{S. Diop\thanks{Charg\'e de Recherche CNRS - LAGEP, UCB Lyon I - 43 Bd.
11 Nov. 1918, 69622 Villeurbanne France;
work supported in part by a grant from NSF/CNRS and in part 
by the Control Systems Laboratory of
the Dept. Elect. Eng. and Comp. Sci., U. of Michigan.}, 
J.\ W.\ Grizzle\thanks{Control Systems Laboratory, \
Department of Electrical Engineering and Computer
Science, \ University of Michigan, Ann Arbor, MI 48109-2122;
work supported in part by the National Science
Foundation under contract NSF ECS-92-13551; matching funds to this grant were
provided by FORD MO.\ CO.},
P.\ E.\ Moraal\footnotemark[2] and A. Stefanopoulou\footnotemark[2] }
%\\[-48pt]}



\maketitle


\begin{abstract}

This paper explores interpolation and numerical differentiation
as a basis for constructing a new approach to the design of nonlinear
observers.  

\end{abstract}


\section{Introduction.}  Observer design is a fundamental problem in system theory and 
control practice.  
The specific problem addressed herein concerns the situation where
one seeks to ``estimate'' the states of a continuous-time model from
observations collected at discrete instants in time.  To date, this problem has
been addressed only in the context of sampled data systems
\cite{grizzle1990b,moraal1992}, where the
procedure has been to first compute
a discrete-time model from the continuous-time model, and
then to design an observer on the basis of the discretized system model.
Of course, the discretization process sometimes can be quite non-trivial, from
a numerical or computational point of view; also, it is usually quite important 
to assume that the observed data has been collected at regular intervals of time,
and thus dealing with problems where event driven sampling is necessary can be difficult.  

This paper arises from the observation that
on the one hand, the numerical analysis literature contains a
substantial body of results on numerical differentiation, while on
the other hand, one of the very basic problems in system theory turns
out to be that of obtaining the time derivatives of a continuous system
variable, which is usually known only through the differential equations of
the system, and time samples of this variable.
We present some initial results on applying numerical differentiation
to the observer design problem; this is done mainly by
working out some specific examples, though the presentation is done
in such a way as to suggest a number of formal ``convergence results'' which will be
investigated in a future publication.

To see how numerical differentiation and observer design may tie together, 
consider the 
following simple observer design problem.
Let a system be given by
\begin{equation}
\left\{
\begin{array}{rcl}
\dot x_1 & = & x_2\\
\dot x_2 & = & -ux_1^3\\
y & = & x_1\end{array}
\right. \label{1}
\end{equation}
and let the question be to build an observer for (\ref{1}); that is,
to reconstruct the states of the system on the basis of the available
measurements.
It is clear that the variable $x_1$ need not be computed since it 
is directly given by the measurement: $ x_1 = y$.
From the equations of (\ref{1}), we see that $x_2 = \dot y$.
If, as is usually the case, we have no access to 
$\dot y$ then we need to {\it compute} $\dot y$ from the measurement $y$.
This is where numerical differentiation comes into play:
we explore the feasibility of {\it numerically
differentiating} $y$ in order to obtain $x_2$.
If $y$ is known only through its time samples, then numerical 
differentiation yields an approximation of
$\dot y$ at the current sampling time $t_k$ by some number
$\hat {\dot y}(t_k)$.
An algorithm for the computation of $\hat {\dot y}(t_k)$ from the
samples of $y$, and {\it estimates of the error bounds} should thus be worked
out \footnote{Ax extensive review of the numerical differntiation and
interpolation
literature is given in \cite{UMTechReport93}.}.
Usually, not only $\dot y$, but, also, finitely many higher
derivatives of $y$ should be approximated at time $t_k$ in order to be
able to compute all of the desired state components.

To our knowledge, the work most closely related to that presented here
is \cite{Sch73}.

\section{An observability property.}\setcounter{equation}{0}
\label{observability}

Given a system described by the state equations
\begin{equation}
\left\{
\begin{array}{rcl}
\dot x_i & = g_i(x, u) \quad (1 \le i \le n)\\
y_j & = h_j(x, u) \quad (1 \le j \le p)\end{array}
\right.
\label{2}
\end{equation}
we will adopt the following specific observability property (see
\cite{grizzle1990b,moraal1992}):
system (\ref{2}) is {\bf observable} if there is an integer $N$ such
that the map $H\left(u, \dot u, \cdots, u^{(N-2)}, \cdot \right)$ of
the state space into some space of output values, defined by
\begin{equation}
H\left(u, \dot u, \cdots, u^{(N-1)}, x \right) = 
\left( 
\begin{array}{c} y \\ \dot
y \\ \vdots \\ y^{(N-1)} \end{array}\right), \label{3}
\end{equation}
is {\it injective} for any fixed {\it universal input}.
It is easiest at this stage to assume that there is actually
no input in (\ref{2}), or that all inputs are universal in the sense that
the above map is injective for any input. 

For an observable system (\ref{2}), we then may write
$$x = L\left(u, \dot u, \cdots, u^{(N-2)}, y, \dot y, \cdots, y^{(N-1)} 
\right).$$
The existence of the map $L$ is guaranteed by our definition of 
observability, but an explicit expression for $L$ may not be easy to 
obtain so that numerical equation solving methods, such as Newton's algorithm,
may be required in order to obtain $L$.  
When (\ref{2}) is a rational system in the sense that the $g_i$'s and the 
$h_j$'s are rational functions of their arguments, then the type of 
observability we just defined corresponds to {\it rational 
observability} of \cite{diop1991d,diop1992b}, and then the explicit
expression of $L$ may be  constructively derived through the use of certain
differential
algebraic techniques.

For observable systems, the observer design problem may thus be 
seen as a problem of numerical differentiation: once estimates of the
derivatives
of $y$ and $u$ can be determined from the available measurements, then $x$ can be
determined from $H$ or $L$.


\section{Illustrative examples.}\setcounter{equation}{0}
\label{examples}

In order to better understand some of the issues involved in the use of  
numerical differentiation and interpolation as a basis for an
observer design method, three academic
examples have been studied.  The first example
involves the estimation
of a time function and a few of its derivatives from noisy samples collected at
discrete instants of time.  The second example builds upon the first one
by showing that if the signal is generated by a known,
continuous-time linear model,
then the model's dynamcis can be incorporated into the estimation
process with advantage.  This will bring us into contact with more classical observer
theory.  The third example was selected to illustrate that the techniques sketched in the
first two examples are
also applicable to nonlinear system models, which is of course, the main
point.  Due to space limitations, it cannot be presented here; the interested
reader is referred to \cite{UMTechReport93}. 


In each case, polynomials are used as the interpolating
funtions, total least squares at the interpolating points is the 
metric to be minimized, and the discrete data will be gathered at a uniform rate.
These simplifications are made so that the discussion can be focused on other issues.

It is emphasized that no theorems are formulated and thus no proofs of validity
are
offered for any of the ``algorithms'' proposed herein.  This paper is meant to be an
exploration of a set of ideas.  Formalization of these ideas into a rigorous design
methodology will be pursued in a later publication.  

\subsection{Sum of two sinusoids.}

The purpose of this example is to illustrate the use of numerical
differentiation, alone and 
in combination with standard system theoretic notions,
to recover or estimate several derivatives of a signal in noise.  The key
parameters of interest are: $N$, the order of the interpolating polynomial;
$\Delta t$, the time interval between data points;
$W * \Delta t$, the length in time of the moving window used for data interpolation
(note that $W + 1$ is the number of data points in the window);
$K$, the data node or ``knot'' within the moving window 
where the derivative is to be estimated (counted from the left, with
the first node in the window numbered zero); and $\sigma$, the ``intensity'' 
of the noise process corrupting the measurements.

Let an analog signal be given by the sum of two sinusoids of frequencies 1 Hz. and
5 Hz., respectively:
\begin{equation}
y(t) = sin(2 \pi t) + sin(10 \pi t).
\label{sines}
\end{equation}

\begin{figure}[hbt]
\centerline{\psfig{file=sumofsines.eps,width=3.5in}}
\smallskip
\small
\caption{\small Time plots of $y(t)$, $y(t) + w(t)$, $y^m_k$ and}
$y^m_k + w_k$, respectively.
\label{signals}
\end{figure}


Assume that the measured signal, $y^m$, is the sum of $y$ and $w$, where $w$ is a
Gaussian white noise process, with standard deviation $\sigma$.  
It is further supposed that
$y^m$ is to be sampled at discrete instants of time,
$k \Delta t$, $k = 0, 1, ...$.  In the absence of noise, the minimum sampling rate 
to reconstruct $y$ from sampled data would be the Nyquist rate, 10 Hz.; with
noise,
at least four times this rate is warranted (an anti-aliasing filter is not
being used
though it would be standard practice to do so).
Discrete samples of $y + w$ will be collected therefore
at 40 Hz.:
\begin{eqnarray}
y^m(t) &:=& y(t) + w(t) \\
y^m_k & :=& y^m(k \Delta t). \nonumber
\label{ysampled}
\end{eqnarray}
%where $G_{BW}$ is the Butterworth, anti-aliasing filter.  
The time signals discussed are depicted in Figure \ref{signals}.  The noise
intensity was selected to provide a (numerically computed)
signal to noise ratio of 20/1, and
corresponded to $\sigma = 0.22$. This should 
provide a feel for the sample rate and noise intensity being used in all later
simulations.


Let the interpolating polynomial for the window of 
data $ \{ y_{k-W}^m, \ldots, y_k^m \}$
be denoted by
\begin{equation}
\hat{y}_k(t) = a_0 + a_1 (t - t_{k-W}) + \ldots + a_N (t - t_{k-W})^N,
\label{interppoly}
\end{equation}
where $t_k := k \Delta t$.
The coefficients $\{ a_0, a_1, \ldots, a_N \}$ are determined from the least
squares solution of 
\begin{eqnarray}
\left[ \begin{array}{cccc} 1 & 0 & \ldots & 0 \\
1 & \Delta t & \ldots & (\Delta t)^N \\
\vdots & \vdots & \vdots & \vdots \\
1 & W \Delta t & \ldots & (W \Delta t)^N 
\end{array} \right] \left[ \begin{array}{l} a_0 \\ a_1 \\ \vdots \\ a_N
\end{array} \right] & = & \left[ \begin{array}{l} y^m_{k-W} \\
y^m_{k-W+1} \\ \vdots \\ y^m_k
\end{array} \right],
\label{leastsquares}
\end{eqnarray}
with respect to the Euclidean norm.  
The estimates of the derivatives of $y$ at time $\bar{t}$ are determined by
\begin{equation}
\widehat{ \frac{d^j}{dt^j}y_k(\bar{t}) } := \frac{d^j}{dt^j} \hat{y}_k(t)_{|t=\bar{t}} \  ;
\label{derivestimates}
\end{equation}
for simplicity of notation, this is written as 
$\frac{d^j}{dt^j} \hat{y}_k( \bar{t} )$.


A number of numerical experiments or simulations
were conducted to ascertain the effects of the
parameters $N$,  $\Delta t$, $W$, $K$ and $\sigma$ on the ability to 
estimate the derivatives
of $y$. 
The general trends of these experiments are summarized here; quantifying
these observations analytically would be a worthy goal:
\begin{itemize}

\item The order of the interpolating polynomial should be selected as low as 
possible in order to average out the noise in the signal. 
On the other hand, for the estimation of 
the 3rd derivative of a signal, for example, 
at least a 3rd order polynomial is required.

\item The window length $W \Delta t$ must be chosen large enough to
capture the variations of the signal so that the interpolating polymomial 
can reproduce its derivatives.  However, larger $W \Delta t$ require
higher order polynomials to be chosen.

\item Estimates of the derivatives made at the end-points of the window
$W \Delta t$ are considerably less accurate than those made in the interior.
However, selecting $K < W$ introduces a time delay of $(W-K) \Delta t$
in the estimates of the signal and its derivatives. 

\item For fixed $N$ and $W \Delta t$, decreasing the time between samples,
$\Delta t$, or increasing
the number of points in the window, $W$, allows a higher noise intensity
$\sigma$ to be tolerated.

\item For a fixed $\Delta t$, though increasing $W$ can produce more accurate results,
it must be considered that it requires more on line compuation and a larger
delay must be introduced (i.e., $(W-K)\Delta t$ is increased)
in order to retain the accuracy of the estimates. 

\end{itemize}

\begin{figure}[hbt]
\centerline{\psfig{file=mplot11.eps,width=3.5in}}
\smallskip
\small
\caption{\small Comparison of the true and estimated \ $(\times)$ val-} 
ues of $y, \ \frac{dy}{dt}, \ 
\frac{d^2y}{dt^2}, 
 \mbox{and} \
\frac{d^3y}{dt^3}$,  respectively using interpolation/numerical differentiation.
\label{figmplot10}
\end{figure}

Figure \ref{figmplot10} compares the true and estimated values of
$y, \ \frac{dy}{dt}, \
\frac{d^2y}{dt^2}, \mbox{and} \
\frac{d^3y}{dt^3}$ when $N=4$, $\Delta t = .025$, $W=8$ and $K=5$. 
Figure \ref{figsineinterp} displays the normalized errors
in the estimates: the errors in the estimates 
for $(y, \ldots, \frac{d^3}{dt^3}y)$ are normalized
by 
\begin{equation}
lim_{T \rightarrow \infty} \frac{1}{T} \int_{0}^{T}{|\frac{d^j}{dt^j}y(t)| dt}.
\label{scales}
\end{equation}
Thus, the error in $y$ is divided by 0.796, $\frac{dy}{dt}$ by 19.3,
$\frac{d^2y}{dt^2}$ by 650, 
and $\frac{d^3y}{dt^3}$ by 22,700.  The wide range in the magnitudes of y
and its derivatives makes the
estimation problem quite challenging.


\begin{figure}[hbt]
\centerline{\psfig{file=mplot9.eps,width=3.5in}}
%\centerline{\psfig{file=sine_interp.eps}}
\smallskip
\small
\caption{\small Normalized estimation errors in $y, \ \frac{dy}{dt}, \
\frac{d^2y}{dt^2}$} 
$\mbox{and} \
\frac{d^3y}{dt^3}$, respectively using interpolation/numerical differentiation.
\label{figsineinterp}
\end{figure}


Estimating the signal and its derivatives on the basis of 
the $W + 1$ data points in the moving window corresponds to using
an FIR (finite impulse response) 
filter.  By introducing a simple modification
to the above, a sort of IIR (infinite
impulse response) filter can be achieved.  The modification is somewhat
analogous to 
what is done in spline
approximations, but to a systems person, it would be called a state.

The idea is to append to (\ref{leastsquares}) the estimates of $y$ and
its derivatives at node K from the previous window, that is, one appends
the relations
\begin{equation}
\frac{d^j}{dt^j} \hat{y}_k(t_{k-1-W+K}) =  \frac{d^j}{dt^j} \hat{y}_{k-1}(t_{k-1-W+K}),
\label{derivconstraints}
\end{equation}
for $j=0, \ldots, J$, for some $0 \le J \le N$.
This is equivalent to
\begin{eqnarray} {\tiny
\left[ \begin{array}{ccccccc} 
1 & (K-1) \Delta t & ((K-1)\Delta t)^2 & \ldots & ((K-1) \Delta t)^J & 
\ldots & ((K-1) \Delta t)^N \\
0 & 1 & 2(K-1) \Delta t & \ldots & J((K-1) \Delta t)^{J-1} & \ldots & 
N((K-1) \Delta t)^{N-1} \\
\vdots & \vdots & \vdots & \vdots & \vdots & \vdots & \vdots \\
0 & 0 & \ldots & \ldots & J! & \ldots & N(N-1)\ldots(N-J+1)((K-1) \Delta t)^{N-J}
\end{array} \right] } \times & \mbox{ } &  \nonumber \\
\left[ \begin{array}{l} a_0 \\ a_1 \\ \vdots \\ a_N \end{array} \right] 
 =  \left[ \begin{array}{r} 
\hat{y}_{k-1}(t_{k-1-W+K}) \\
\frac{d}{dt} \hat{y}_{k-1}(t_{k-1-W+K}) \\ 
\vdots \\ 
\frac{d^J}{dt^J} \hat{y}_{k-1}(t_{k-1-W+K})
\end{array} \right] & \mbox{\hspace*{.2in}} &.
\label{matrixderivconstraints}
\end{eqnarray}

The right-hand side of (\ref{derivconstraints}) must be initialized to
start the interpolation process, and thus becomes a state.  This state 
allows derivative 
information to be passed from one window of data 
to the next.   
Weights can be added to trade-off errors in interpolating the measured data points
versus errors in interpolating the derivatives in (\ref{derivconstraints}). 
Due to space limitations, no simulations employing this
modification are presented; a related idea is explored in the next subsection.

\subsection{Including a system model in the interpolation process.}

The goal of this example is to demonstrate two methods for incorporating
a dynamic model of a system into the interpolation/numerical differentiation
process,
whenever the signal in question is generated by a finite set of ordinary
differential equations.
This will bring us into more direct contact with observer design because we
will tightly connect estimating the derivatives of a system's output
with estimating the states of the model.

The signal (\ref{sines}) can obviously be generated by an initialized,
uncontrolled, linear model with four states.  Let 
\begin{equation}
\begin{array}{rcl}
\frac{dx}{dt} & = &  A x  \\
y & = & C x 
\end{array}
\label{linearmodel}
\end{equation}
be such a model; it will be observable.  This means that $x$ can be recovered
from estimates of $y$ and its derivatives through
\begin{equation} \left[ \begin{array}{r} 
y \\ \frac{d}{dt}y \\ \vdots \\ \frac{d^{\eta}}{dt^{\eta}}y
\end{array} \right] = \left[ \begin{array}{l} 
C \\ CA \\ \vdots \\ CA^{\eta}
\end{array} \right] x
\label{obsmat}
\end{equation}
for $\eta$ large enough, which in the particular case of (\ref{linearmodel}) is
$\eta = 3$.

The first way to include the model in the estimation process is now
explained.  Suppose that in general the model (\ref{linearmodel}) has dimension
$n$, and for simplicity, assume that it has a single output.
As before, let
\begin{equation}
\left[ \begin{array}{r} 
\hat{y}_{k-1}(t_{k-1-W+K}) \\
\frac{d}{dt} \hat{y}_{k-1}(t_{k-1-W+K}) \\ 
\vdots \\ 
\frac{d^{n-1}}{dt^{n-1}} \hat{y}_{k-1}(t_{k-1-W+K})
\end{array} \right]
\label{ydotest}
\end{equation}
be estimates of $y$ and its first $n-1$ derivatives from the previous window.
By substituting (\ref{ydotest}) for the right-hand side of (\ref{obsmat}) with
$\eta = n-1$, one can
determine $\hat{x}_{k-1}(t_{k-1-W+K})$. 
If the resulting estimate of $x$ is ``accurate'', then it satisfies (\ref{obsmat}) for
any $\eta$; if not, choosing $\eta \ge n$ introduces constraints whose
errors can be minimized in order for $x$ to be more closely compatible
with the dynamics (\ref{linearmodel}).  Thus, in place of
(\ref{derivconstraints}),
one augments (\ref{leastsquares}) with
\begin{equation}
\frac{d^j}{dt^j} \hat{y}_k(t_{k-1-W+K}) = CA^j  \hat{x}_{k-1}(t_{k-1-W+K}),
\label{modelconstraints}
\end{equation}
For $j=0, \ldots, \eta$, and proceeds as indicated previously.  
Note that taking $\eta =J \le n-1$ reduces 
(\ref{modelconstraints}) to (\ref{derivconstraints}).



\begin{figure}[hbt]
\centerline{\psfig{file=sine_spline_order5.eps,width=3.5in}}
\smallskip
\caption{\small Normalized estimation errors in $y, \ \frac{dy}{dt}, \
\frac{d^2y}{dt^2}$,}
$ \mbox{and} \
\frac{d^3y}{dt^3}$, respectively using interpolation/numerical differentiation
and $\eta=5$ in (\protect\ref{modelconstraints}).  Note the reduced error with respect to
Figure \protect\ref{figsineinterp}. 
\label{figsinesplineorder5}
\end{figure}


Figure \ref{figsinesplineorder5} shows the
effect of this modification on the estimation process for $N=6$, $\Delta t = .025$, 
$W=8$ and $K=5$.  The weight on the constraints was chosen as $ Q =$
\begin{equation}
\left[ \begin{array}{cccccc} 1.26 & 0 & 0 & 0 & 0 & 0\\
0  & 5.19 \ 10^{-2}  & 0 & 0 & 0 & 0 \\
0 & 0 & 1.54 \ 10^{-3} & 0 & 0 & 0 \\
0 & 0 & 0 & 4.4 \ 10^{-5} & 0 & 0 \\
0 & 0 & 0 & 0 & 1.3 \ 10^{-6} & 0 \\
0 & 0 & 0 & 0 & 0 & 4.4 \ 10^{-8}
\end{array} \right].
\end{equation}
The rationale for
selecting the entries of the
weight $Q$ is that the magnitude of each successive derivative grows by a
approximatley
30 (recall (\ref{scales}), or see (\ref{sines})), 
and thus $Q$ corresponds to equal ``relative
weighting'' on $y$ and its derivatives;
this same rationale will be applied to subsequent examples. Including the model
interpolation constraints allowed the same window size and delay as in Figure
\ref{figsineinterp} to be used with an increased polynomial order, without
``following'' the noise.


The method just outlined for estimating $x$ from $y$ and its derivatives is akin to
a ``dead-beat'' observer for $x$, though, as explained
earlier, there is smoothing in the
interpolation of $y$, which is ``IIR''.  Some additional smoothing from the
model can be obtained, and this constitutes the second method for introducing
the model into the interpolation/numerical differentiation process.
This technique, which is reminiscent of
\cite{grizzle1990b,moraal1992}, 
will be
directly applicable to nonlinear systems.  

Let $A_d$ denote the discretization of the linear model, (\ref{linearmodel}),
and let $\hat{x}_{k-1}^{+}(t_{k-1-W+K})$ be the previous estimate of $x$.  Update $x$
based on the latest estimate of $y$ and its derivatives through a damped Newton
method:
\begin{eqnarray}
\hat{x}_{k}^{-}(t_{k-W+K}) & = & A_d \hat{x}_{k-1}^{+}(t_{k-1-W+K}) \nonumber \\
\hat{x}_{k}^{+}(t_{k-W+K}) & = & \hat{x}_{k}^{-}(t_{k-W+K}) + \label{newton} \\
\epsilon \left[ \begin{array}{l} C \\ CA \\ \vdots \\ CA^{n-1} \end{array} 
\right]^{-1} & &  ( \left[ \begin{array}{l} 
\hat{y}_{k}(t_{k-W+K}) \\
\frac{d}{dt} \hat{y}_{k}(t_{k-W+K}) \\ 
\vdots \\ 
\frac{d^{n-1}}{dt^{n-1}} \hat{y}_{k}(t_{k-W+K}) \end{array} \right] - 
\left[ \begin{array}{l} C \\ CA \\ \vdots \\ CA^{n-1} \end{array} 
\right]  \hat{x}_{k}^{-}(t_{k-W+K}) ), \nonumber
\end{eqnarray}
where $0 < \epsilon <  1$; using $\epsilon = 1$ reduces the above to a ``dead-beat''
estimator for $x$.  The damped Newton update can be used with
(\ref{leastsquares}) alone or with any of
the additions to (\ref{leastsquares}) discussed previously.

Selected simulations of the above estimation methods are now presented.
In what follows,
the state in (\ref{linearmodel}) was chosen as 
$$ x = \left[ \begin{array}{r}
y \\ \frac{d}{dt} y \\  \frac{d^2}{dt^2} y \\ \frac{d^3}{dt^3} y 
\end{array} \right]. $$  
It was properly initialized so that its output is precisely
given by (\ref{sines}).  The same noise sequence used in
Figure \ref{figsineinterp} was added to the output of (\ref{linearmodel}).
%also, the anti-aliasing filter (\ref{ysampled}) was used.


\begin{figure}[hbt]
\centerline{\psfig{file=sine_spline.eps,width=3.5in}}

\caption{\small Normalized estimation errors in \protect$y, \ \frac{dy}{dt}, \
\frac{d^2y}{dt^2}$} 
and  \protect$\frac{d^3y}{dt^3}$, respectively, incorporating 
``spline method'' and
Newton update in interpolation/numerical differentiation.
\label{figsinespline}
\end{figure}


Figure \ref{figsinespline} shows the normalized estimation error resulting
from the application of (\ref{leastsquares})  in combination with
(\ref{derivconstraints})
and (\ref{newton}).   The parameter values were chosen as: $N=5$, $\Delta t = 0.025$
sec., $W=6$ and $K=3$  in (\ref{leastsquares}), $\eta = 4$ in 
(\ref{modelconstraints}) with
a weight on (\ref{derivconstraints}) selected as
\begin{equation} Q = 0.1 \times
\left[ \begin{array}{cccc} 1.26 \, 10^{-1} & 0 & 0 & 0 \\
0  & 5.19 \, 10^{-3}  & 0 & 0 \\
0 & 0 & 1.54 \, 10^{-4} & 0 \\
0 & 0 & 0 & 4.4 \, 10^{-5} 
\end{array} \right],
\end{equation}
and $\epsilon = 0.03$ in (\ref{newton}).   
For comparison purposes,
a steady state, discrete-time, Kalman filter was designed with 
state noise covariance equal to 
$0.01 \ \sigma^2 \ Q^{-1}$ and output noise variance equal to $\sigma^2$; even
though (\ref{linearmodel}) does not have any process noise per se, the pair
consisting of the matrix $A$ and the state noise convariance matrix
had to be made stabilizable
for a stable Kalman filter to exist.
The magnitudes of the steady state errors were essentially identical to those in 
Figure \ref{figsinespline}.


\begin{figure}[hbt]
\centerline{\psfig{file=sine_model.eps,width=3.5in}}
\caption{\small Normalized estimation errors in $y, \ \frac{dy}{dt}, \
\frac{d^2y}{dt^2} \ $} 
$\mbox{and} \ \frac{d^3y}{dt^3}$, respectively, incorporating model derivative
constraints and damped Newton updates with 
interpolation/numerical differentiation.
\label{figsinemodel}
\end{figure}


The normalized estimation errors resulting
from the application of (\ref{leastsquares}) in combination with
(\ref{modelconstraints})
and (\ref{newton}) are shown in
Figure \ref{figsinemodel}.   The parameter values were chosen as: $N=5$, 
$\Delta t = 0.025$
sec., $W=6$ and $K=3$  in (\ref{leastsquares}), $\eta = 4$ in
(\ref{modelconstraints}) 
with a weight selected as
\begin{equation} Q = 
\left[ \begin{array}{ccccc} 1.26 \, 10^{-1} & 0 & 0 & 0 & 0 \\
0  & 5.19 \, 10^{-3}  & 0 & 0 & 0 \\
0 & 0 & 1.54 \, 10^{-4} & 0 & 0 \\
0 & 0 & 0 & 4.4 \, 10^{-5} & 0 \\
0 & 0 & 0 & 0 & 1.33 \, 10^{-6}
\end{array} \right];
\end{equation}

\begin{thebibliography}{11}


\bibitem{diop1992b}
S.~Diop.
\newblock Rational system equivalence, and generalized realization theory.
\newblock In {\em Proceedings of the {I}FAC-{S}ymposium {N}OLCOS'92}, pages
  402--407, Bordeaux, 1992. IFAC.

\bibitem{diop1991d}
S.~Diop and M.~Fliess.
\newblock On nonlinear observability.
\newblock In C.~Commault, D.~Normand-Cyrot, J.~M. Dion, L.~Dugard, M.~Fliess,
  A.~Titli, G.~Cohen, A.~Benveniste, and I.~D. Landau, editors, {\em
  Proceedings of the {F}irst {E}uropean {C}ontrol {C}onference}, pages
  152--157, Paris, 1991. Herm{\`e}s.

\bibitem{UMTechReport93}
S. Diop, J.\ W.\ Grizzle, 
P.\ E.\ Moraal and A. Stefanopoulou, ``Interpolation and Numerical 
Differentiation for Observer Design'', University Michigan's College of
Engineering Control Group Report Series,  Report No. CGR-93-14, September
1993.


\bibitem{grizzle1990b}
J.~W. Grizzle and P.~E. Moraal.
\newblock Newton observers and nonlinear discrete-time control.
\newblock In {\em Proceedings of the 29th {I}EEE {C}onference on {D}ecision and
  {C}ontrol}, 1990, pages 760--767.

\bibitem{moraal1992}
P.~E. Moraal and J.~W. Grizzle.
\newblock Nonlinear discrete-time observers using {N}ewton's and {B}royden's
  methods.
\newblock In {\em Proceedings of the {A}merical {C}ontrol {C}onference}, June
1992, pages 3086--3090.

\bibitem{Sch73}
U. Sch\"{o}nwandt, ``Approximations to nonlinear observers'', {\it Automatica},
Vol. 9,1973,  pages 349--356.

\end{thebibliography}

\end{document}



From moraal@eecs.umich.edu Fri Jan 21 10:46:18 1994
Date: Thu, 20 Jan 1994 17:35:08 -0500
From: Paul E Moraal <moraal@eecs.umich.edu>
To: grizzle@eecs.umich.edu
Subject: numdiff_ACC.tex

\documentstyle[psfig,proc,sectspace]{article}
\pagestyle{empty}
%\input define
\textheight=24cm
\textwidth=17.7cm
\oddsidemargin=-.4cm
\columnsep=1.5cm
\topmargin=-2.716cm
%\def\thesection{\arabic{section}.}
%\def\thesubsection{\thesection\arabic{subsection}.}
%\def\theequation{\arabic{equation}}
%\def\figurename{\small Figure}
%******************************************************
%******From numdiff_SEPT8.tex

\newcommand{\r}{\hbox{\sl l\kern-.2em R}}
\def\thesection{\arabic{section}}
\def\thesubsection{\thesection.\arabic{subsection}}
\def\theequation{\thesection.\arabic{equation}}
\def\theenumi{\roman{enumi}}\def\labelenumi{(\theenumi).}
\def\theenumii{\alph{enumii}}\def\labelenumii{(\theenumii)}
\newtheorem{theorem}{Theorem}[section]
\newtheorem{proof}{Proof:}
\def\theproof{}
\def\thetheorem{\thesection\arabic{theorem}}

\def\setR{\hbox{$l\! \! R$}}
\def\setN{\hbox{$l\! \! N$}}
\def\CN{{\cal N}}

\def \BA {{\bf A}}
\def \BB {{\bf B}}
\def \BC {{\bf C}}
\def \BE {{\bf E}}
\def \BF {{\bf F}}
\def \BG {{\bf G}}
\def \BH {{\bf H}}
\def \BI {{\bf I}}
\def \BJ {{\bf J}}
\def \Bk {{\bf k}}
\def \BK {{\bf K}}
\def \BL {{\bf L}}
\def \BM {{\bf M}}
\def \BN {{\bf N}}
\def \BQ {{\bf Q}}
\def \BP {{\bf P}}
\def \BQ {{\bf Q}}
\def \BR {{\bf R}}
\def \BS {{\bf S}}
\def \BV {{\bf V}}

%\def \CN   {{\eusm N}}

\def\thesubsubsection{\thesubsection\arabic{subsubsection}.}

%******************************************************

\makeatletter
\def\section{\@startsection{section}{1}{\z@}{\sectionspaceb
plus .5em minus .1em}{\sectionspacea plus 0em minus
.1em}{\centering\normalsize\bf}}

\def\subsection{\@startsection{subsection}{2}{\z@}{\sectionspaceb plus .5em minus
.1em}{\sectionspacea plus 0em minus .1em}{\normalsize\bf}}

\def\subsubsection{\@startsection{subsubsection}{3}{\z@}{\sectionspaceb plus .5em
minus .1em}{\qquad}{\underline{\normalsize\bf}}}

\def\caption{\refstepcounter\@captype \@dblarg{\@caption\@captype}}

%\long\def\@makecaption#1#2{\vskip 10pt \setbox\@tempboxa\hbox{#1: #2}
%\ifdim \wd\@tempboxa >\hsize #1: #2\par \else \hbox
%to \hsize{\hfil\box\@tempboxa\hfil} \fi}

\long\def\@makecaption#1#2{\hbox{#1: #2}}

\long\def\@caption#1[#2]#3{\addcontentsline{\csname
  ext@#1\endcsname}{#1}{\protect\numberline{\csname 
  the#1\endcsname}{\ignorespaces #2}}\par
  \begingroup
    \@parboxrestore
    \normalsize
    \@makecaption{\csname fnum@#1\endcsname}{\ignorespaces #3}\par
  \endgroup}

\makeatother

\begin{document}

\small

\baselineskip=12pt
\parskip=0in


\title{\vspace*{.5cm}
\large\bf Interpolation and Numerical Differentiation \\
for Observer Design}

\author{S. Diop\thanks{Charg\'e de Recherche CNRS - LAGEP, UCB Lyon I - 43 Bd.
11 Nov. 1918, 69622 Villeurbanne France;
work supported in part by a grant from NSF/CNRS and in part 
by the Control Systems Laboratory of
the Dept. Elect. Eng. and Comp. Sci., U. of Michigan.}, 
J.\ W.\ Grizzle\thanks{Control Systems Laboratory, \
Department of Electrical Engineering and Computer
Science, \ University of Michigan, Ann Arbor, MI 48109-2122;
work supported in part by the National Science
Foundation under contract NSF ECS-92-13551; matching funds to this grant were
provided by FORD MO.\ CO.},
P.\ E.\ Moraal\footnotemark[2] and A. Stefanopoulou\footnotemark[2] }
%\\[-48pt]}



\maketitle


\begin{abstract}

This paper explores interpolation and numerical differentiation
as a basis for constructing a new approach to the design of nonlinear
observers.  

\end{abstract}


\section{Introduction.}  Observer design is a fundamental problem in system theory and 
control practice.  
The specific problem addressed herein concerns the situation where
one seeks to ``estimate'' the states of a continuous-time model from
observations collected at discrete instants in time.  To date, this problem has
been addressed only in the context of sampled data systems
\cite{grizzle1990b,moraal1992}, where the
procedure has been to first compute
a discrete-time model from the continuous-time model, and
then to design an observer on the basis of the discretized system model.
Of course, the discretization process sometimes can be quite non-trivial, from
a numerical or computational point of view; also, it is usually quite important 
to assume that the observed data has been collected at regular intervals of time,
and thus dealing with problems where event driven sampling is necessary can be difficult.  

This paper arises from the observation that
on the one hand, the numerical analysis literature contains a
substantial body of results on numerical differentiation, while on
the other hand, one of the very basic problems in system theory turns
out to be that of obtaining the time derivatives of a continuous system
variable, which is usually known only through the differential equations of
the system, and time samples of this variable.
We present some initial results on applying numerical differentiation
to the observer design problem; this is done mainly by
working out some specific examples, though the presentation is done
in such a way as to suggest a number of formal ``convergence results'' which will be
investigated in a future publication.

To see how numerical differentiation and observer design may tie together, 
consider the 
following simple observer design problem.
Let a system be given by
\begin{equation}
\left\{
\begin{array}{rcl}
\dot x_1 & = & x_2\\
\dot x_2 & = & -ux_1^3\\
y & = & x_1\end{array}
\right. \label{1}
\end{equation}
and let the question be to build an observer for (\ref{1}); that is,
to reconstruct the states of the system on the basis of the available
measurements.
It is clear that the variable $x_1$ need not be computed since it 
is directly given by the measurement: $ x_1 = y$.
>From the equations of (\ref{1}), we see that $x_2 = \dot y$.
If, as is usually the case, we have no access to 
$\dot y$ then we need to {\it compute} $\dot y$ from the measurement $y$.
This is where numerical differentiation comes into play:
we explore the feasibility of {\it numerically
differentiating} $y$ in order to obtain $x_2$.
If $y$ is known only through its time samples, then numerical 
differentiation yields an approximation of
$\dot y$ at the current sampling time $t_k$ by some number
$\hat {\dot y}(t_k)$.
An algorithm for the computation of $\hat {\dot y}(t_k)$ from the
samples of $y$, and {\it estimates of the error bounds} should thus be worked
out \footnote{An extensive review of the numerical differentiation and
interpolation
literature is given in \cite{UMTechReport93}.}.
Usually, not only $\dot y$, but, also, finitely many higher
derivatives of $y$ should be approximated at time $t_k$ in order to be
able to compute all of the desired state components.

To our knowledge, the work most closely related to that presented here
is \cite{Sch73}.

\section{An observability property.}\setcounter{equation}{0}
\label{observability}

Given a system described by the state equations
\begin{equation}
\left\{
\begin{array}{rcl}
\dot x_i & = g_i(x, u) \quad (1 \le i \le n)\\
y_j & = h_j(x, u) \quad (1 \le j \le p)\end{array}
\right.
\label{2}
\end{equation}
we will adopt the following specific observability property (see
\cite{grizzle1990b,moraal1992}):
system (\ref{2}) is {\bf observable} if there is an integer $N$ such
that the map $H\left(u, \dot u, \cdots, u^{(N-1)}, \cdot \right)$ of
the state space into some space of output values, defined by
\begin{equation}
H\left(u, \dot u, \cdots, u^{(N-1)}, x \right) = 
\left( 
\begin{array}{c} y \\ \dot
y \\ \vdots \\ y^{(N-1)} \end{array}\right), \label{3}
\end{equation}
is {\it injective} for any fixed {\it universal input}.
It is easiest at this stage to assume that there is actually
no input in (\ref{2}), or that all inputs are universal in the sense that
the above map is injective for any input. 

For an observable system (\ref{2}), we then may write
$$x = L\left(u, \dot u, \cdots, u^{(N-1)}, y, \dot y, \cdots, y^{(N-1)} 
\right).$$
The existence of the map $L$ is guaranteed by our definition of 
observability, but an explicit expression for $L$ may not be easy to 
obtain so that numerical equation solving methods, such as Newton's algorithm,
may be required in order to obtain $L$.  
When (\ref{2}) is a rational system in the sense that the $g_i$'s and the 
$h_j$'s are rational functions of their arguments, then the type of 
observability we just defined corresponds to {\it rational 
observability} of \cite{diop1991d,diop1992b}, and then the explicit
expression of $L$ may be  constructively derived through the use of certain
differential
algebraic techniques.

For observable systems, the observer design problem may thus be 
seen as a problem of numerical differentiation: once estimates of the
derivatives
of $y$ and $u$ can be determined from the available measurements, then $x$ can be
determined from $H$ or $L$.


\section{Illustrative examples.}\setcounter{equation}{0}
\label{examples}

In order to better understand some of the issues involved in the use of  
numerical differentiation and interpolation as a basis for an
observer design method, three academic
examples have been studied.  The first example
involves the estimation
of a time function and a few of its derivatives from noisy samples collected at
discrete instants of time.  The second example builds upon the first one
by showing that if the signal is generated by a known,
continuous-time linear model,
then the model's dynamics can be incorporated into the estimation
process with advantage.  This will bring us into contact with more classical observer
theory.  The third example was selected to illustrate that the techniques sketched in the
first two examples are
also applicable to nonlinear system models, which is of course, the main
point.  Due to space limitations, it cannot be presented here; the interested
reader is referred to \cite{UMTechReport93}. 


In each case, polynomials are used as the interpolating
functions, total least squares at the interpolating points is the 
metric to be minimized, and the discrete data will be gathered at a uniform rate.
These simplifications are made so that the discussion can be focused on other issues.

It is emphasized that no theorems are formulated and thus no proofs of validity
are
offered for any of the ``algorithms'' proposed herein.  This paper is meant to be an
exploration of a set of ideas.  Formalization of these ideas into a rigorous design
methodology will be pursued in a later publication.  

\subsection{Sum of two sinusoids.}

The purpose of this example is to illustrate the use of numerical
differentiation, alone and 
in combination with standard system theoretic notions,
to recover or estimate several derivatives of a signal in noise.  The key
parameters of interest are: $N$, the order of the interpolating polynomial;
$\Delta t$, the time interval between data points;
$W * \Delta t$, the length in time of the moving window used for data interpolation
(note that $W + 1$ is the number of data points in the window);
$K$, the data node or ``knot'' within the moving window 
where the derivative is to be estimated (counted from the left, with
the first node in the window numbered zero); and $\sigma$, the ``intensity'' 
of the noise process corrupting the measurements.

Let an analog signal be given by the sum of two sinusoids of frequencies 1 Hz. and
5 Hz., respectively:
\begin{equation}
y(t) = sin(2 \pi t) + sin(10 \pi t).
\label{sines}
\end{equation}

\begin{figure}[hbt]
\centerline{\psfig{file=sumofsines.eps,width=3.5in}}
\smallskip
\small
\caption{\small Time plots of $y(t)$, $y(t) + w(t)$, $y^m_k$ and}
$y^m_k + w_k$, respectively.
\label{signals}
\end{figure}


Assume that the measured signal, $y^m$, is the sum of $y$ and $w$, where $w$ is a
Gaussian white noise process, with standard deviation $\sigma$.  
It is further supposed that
$y^m$ is to be sampled at discrete instants of time,
$k \Delta t$, $k = 0, 1, ...$.  In the absence of noise, the minimum sampling rate 
to reconstruct $y$ from sampled data would be the Nyquist rate, 10 Hz.; with
noise,
at least four times this rate is warranted (an anti-aliasing filter is not
being used
though it would be standard practice to do so).
Discrete samples of $y + w$ will be collected therefore
at 40 Hz.:
\begin{eqnarray}
y^m(t) &:=& y(t) + w(t) \\
y^m_k & :=& y^m(k \Delta t). \nonumber
\label{ysampled}
\end{eqnarray}
%where $G_{BW}$ is the Butterworth, anti-aliasing filter.  
The time signals discussed are depicted in Figure \ref{signals}.  The noise
intensity was selected to provide a (numerically computed)
signal to noise ratio of 20/1, and
corresponded to $\sigma = 0.22$. This should 
provide a feel for the sample rate and noise intensity being used in all later
simulations.


Let the interpolating polynomial for the window of 
data $ \{ y_{k-W}^m, \ldots, y_k^m \}$
be denoted by
\begin{equation}
\hat{y}_k(t) = a_0 + a_1 (t - t_{k-W}) + \ldots + a_N (t - t_{k-W})^N,
\label{interppoly}
\end{equation}
where $t_k := k \Delta t$.
The coefficients $\{ a_0, a_1, \ldots, a_N \}$ are determined from the least
squares solution of 
\begin{equation} \tiny
\left[ \begin{array}{cccc} 1 & 0 & \ldots & 0 \\
1 & \Delta t & \ldots & (\Delta t)^N \\
\vdots & \vdots & \vdots & \vdots \\
1 & W \Delta t & \ldots & (W \Delta t)^N 
\end{array} \right] \left[ \begin{array}{l} a_0 \\ a_1 \\ \vdots \\ a_N
\end{array} \right]  =  \left[ \begin{array}{l} y^m_{k-W} \\
y^m_{k-W+1} \\ \vdots \\ y^m_k
\end{array} \right], \normalsize
\label{leastsquares}
\end{equation}
with respect to the Euclidean norm.  
The estimates of the derivatives of $y$ at time $\bar{t}$ are determined by
\begin{equation}
\widehat{ \frac{d^j}{dt^j}y_k(\bar{t}) } := \frac{d^j}{dt^j} \hat{y}_k(t)_{|t=\bar{t}} \  ;
\label{derivestimates}
\end{equation}
for simplicity of notation, this is written as 
$\frac{d^j}{dt^j} \hat{y}_k( \bar{t} )$.


A number of numerical experiments or simulations
were conducted to ascertain the effects of the
parameters $N$,  $\Delta t$, $W$, $K$ and $\sigma$ on the ability to 
estimate the derivatives
of $y$. 
The general trends of these experiments are summarized here; quantifying
these observations analytically would be a worthy goal:
\begin{itemize}

\item The order of the interpolating polynomial should be selected as low as 
possible in order to average out the noise in the signal. 
On the other hand, for the estimation of 
the 3rd derivative of a signal, for example, 
at least a 3rd order polynomial is required.

\item The window length $W \Delta t$ must be chosen large enough to
capture the variations of the signal so that the interpolating polymomial 
can reproduce its derivatives.  However, larger $W \Delta t$ require
higher order polynomials to be chosen.

\item Estimates of the derivatives made at the end-points of the window
$W \Delta t$ are considerably less accurate than those made in the interior.
However, selecting $K < W$ introduces a time delay of $(W-K) \Delta t$
in the estimates of the signal and its derivatives. 

\item For fixed $N$ and $W \Delta t$, decreasing the time between samples,
$\Delta t$, or increasing
the number of points in the window, $W$, allows a higher noise intensity
$\sigma$ to be tolerated.

\item For a fixed $\Delta t$, though increasing $W$ can produce more accurate results,
it must be considered that it requires more on line compuation and a larger
delay must be introduced (i.e., $(W-K)\Delta t$ is increased)
in order to retain the accuracy of the estimates. 

\end{itemize}

\begin{figure}[hbt]
\centerline{\psfig{file=mplot11.eps,width=3.5in}}
\smallskip
\small
\caption{\small Comparison of the true and estimated \ $(\times)$ val-} 
ues of $y, \ \frac{dy}{dt}, \ 
\frac{d^2y}{dt^2}, 
 \mbox{and} \
\frac{d^3y}{dt^3}$,  respectively using interpolation/numerical differentiation.
\label{figmplot10}
\end{figure}

Figure \ref{figmplot10} compares the true and estimated values of
$y, \ \frac{dy}{dt}, \
\frac{d^2y}{dt^2}, \mbox{and} \
\frac{d^3y}{dt^3}$ when $N=4$, $\Delta t = .025$, $W=8$ and $K=5$. 
Figure \ref{figsineinterp} displays the normalized errors
in the estimates: the errors in the estimates 
for $(y, \ldots, \frac{d^3}{dt^3}y)$ are normalized
by 
\begin{equation}
lim_{T \rightarrow \infty} \frac{1}{T} \int_{0}^{T}{|\frac{d^j}{dt^j}y(t)| dt}.
\label{scales}
\end{equation}
Thus, the error in $y$ is divided by 0.796, $\frac{dy}{dt}$ by 19.3,
$\frac{d^2y}{dt^2}$ by 650, 
and $\frac{d^3y}{dt^3}$ by 22,700.  The wide range in the magnitudes of y
and its derivatives makes the
estimation problem quite challenging.


\begin{figure}[hbt]
\centerline{\psfig{file=mplot9.eps,width=3.5in}}
%\centerline{\psfig{file=sine_interp.eps}}
\smallskip
\small
\caption{\small Normalized estimation errors in $y, \ \frac{dy}{dt}, \
\frac{d^2y}{dt^2}$} 
$\mbox{and} \
\frac{d^3y}{dt^3}$, respectively using interpolation/numerical differentiation 
( eq. \protect (\ref{leastsquares})).
\label{figsineinterp}
\end{figure}


Estimating the signal and its derivatives on the basis of 
the $W + 1$ data points in the moving window corresponds to using
an FIR (finite impulse response) 
filter.  By introducing a simple modification
to the above, a sort of IIR (infinite
impulse response) filter can be achieved.  The modification is somewhat
analogous to 
what is done in spline
approximations, but to a systems person, it would be called a state.

The idea is to append to (\ref{leastsquares}) the estimates of $y$ and
its derivatives at node K from the previous window, that is, one appends
the relations
\begin{equation}
\frac{d^j}{dt^j} \hat{y}_k(t_{k-1-W+K}) =  \frac{d^j}{dt^j} \hat{y}_{k-1}(t_{k-1-W+K}),
\label{derivconstraints}
\end{equation}
for $j=0, \ldots, J$, for some $0 \le J \le N$.
%This is equivalent to
%\begin{eqnarray} {\tiny
%\left[ \begin{array}{ccccccc} 
%1 & (K-1) \Delta t & ((K-1)\Delta t)^2 & \ldots & ((K-1) \Delta t)^J & 
%\ldots & ((K-1) \Delta t)^N \\
%0 & 1 & 2(K-1) \Delta t & \ldots & J((K-1) \Delta t)^{J-1} & \ldots & 
%N((K-1) \Delta t)^{N-1} \\
%\vdots & \vdots & \vdots & \vdots & \vdots & \vdots & \vdots \\
%0 & 0 & \ldots & \ldots & J! & \ldots & N(N-1)\ldots(N-J+1)((K-1) \Delta t)^{N-J}
%\end{array} \right] } \times & \mbox{ } &  \nonumber \\
%\left[ \begin{array}{l} a_0 \\ a_1 \\ \vdots \\ a_N \end{array} \right] 
% =  \left[ \begin{array}{r} 
%\hat{y}_{k-1}(t_{k-1-W+K}) \\
%\frac{d}{dt} \hat{y}_{k-1}(t_{k-1-W+K}) \\ 
%\vdots \\ 
%\frac{d^J}{dt^J} \hat{y}_{k-1}(t_{k-1-W+K})
%\end{array} \right] & \mbox{\hspace*{.2in}} &.
%\label{matrixderivconstraints}
%\end{eqnarray}

The right-hand side of (\ref{derivconstraints}) must be initialized to
start the interpolation process, and thus becomes a state.  This state 
allows derivative 
information to be passed from one window of data 
to the next.   
Weights can be added to trade-off errors in interpolating the measured data points
versus errors in interpolating the derivatives in (\ref{derivconstraints}). 
Due to space limitations, no simulations employing this
modification are presented; a related idea is explored in the next subsection.

\subsection{Including a system model in the interpolation process.}

The goal of this example is to demonstrate two methods for incorporating
a dynamic model of a system into the interpolation/numerical differentiation
process,
whenever the signal in question is generated by a finite set of ordinary
differential equations.
This will bring us into more direct contact with observer design because we
will tightly connect estimating the derivatives of a system's output
with estimating the states of the model.

The signal (\ref{sines}) can obviously be generated by an initialized,
uncontrolled, linear model with four states.  Let 
\begin{equation}
\begin{array}{rcl}
\frac{dx}{dt} & = &  A x  \\
y & = & C x 
\end{array}
\label{linearmodel}
\end{equation}
be such a model; it will be observable.  This means that $x$ can be recovered
>from estimates of $y$ and its derivatives through
\begin{equation} \left[ \begin{array}{r} 
y \\ \frac{d}{dt}y \\ \vdots \\ \frac{d^{\eta}}{dt^{\eta}}y
\end{array} \right] = \left[ \begin{array}{l} 
C \\ CA \\ \vdots \\ CA^{\eta}
\end{array} \right] x
\label{obsmat}
\end{equation}
for $\eta$ large enough, which in the particular case of (\ref{linearmodel}) is
$\eta = 3$.

The first way to include the model in the estimation process is now
explained.  Suppose that in general the model (\ref{linearmodel}) has dimension
$n$, and for simplicity, assume that it has a single output.
As before, let
\begin{equation}
\left[ \begin{array}{c} 
\hat{y}_{k-1}(t_{k-1-W+K}) \\
\frac{d}{dt} \hat{y}_{k-1}(t_{k-1-W+K}) \\ 
\vdots \\ 
\frac{d^{n-1}}{dt^{n-1}} \hat{y}_{k-1}(t_{k-1-W+K})
\end{array} \right]
\label{ydotest}
\end{equation}
be estimates of $y$ and its first $n-1$ derivatives from the previous window.
By substituting (\ref{ydotest}) for the right-hand side of (\ref{obsmat}) with
$\eta = n-1$, one can
determine $\hat{x}_{k-1}(t_{k-1-W+K})$. 
If the resulting estimate of $x$ is ``accurate'', then it satisfies (\ref{obsmat}) for
any $\eta$; if not, choosing $\eta \ge n$ introduces constraints whose
errors can be minimized in order for $x$ to be more closely compatible
with the dynamics (\ref{linearmodel}).  Thus, in place of
(\ref{derivconstraints}),
one augments (\ref{leastsquares}) with
\begin{equation}
\frac{d^j}{dt^j} \hat{y}_k(t_{k-1-W+K}) = CA^j  \hat{x}_{k-1}(t_{k-1-W+K}),
\label{modelconstraints}
\end{equation}
For $j=0, \ldots, \eta$, and proceeds as indicated previously.  
Note that taking $\eta =J \le n-1$ reduces 
(\ref{modelconstraints}) to (\ref{derivconstraints}).



\begin{figure}[hbt]
\centerline{\psfig{file=sine_spline_order5.eps,width=3.5in}}
\smallskip
\caption{\small Normalized estimation errors in $y, \ \frac{dy}{dt}, \
\frac{d^2y}{dt^2}$,}
$ \mbox{and} \
\frac{d^3y}{dt^3}$, respectively using interpolation/numerical differentiation
(\protect\ref{leastsquares}) and $\eta=5$ in (\protect\ref{modelconstraints}).  Note the reduced error with respect to
Figure \protect\ref{figsineinterp}. 
\label{figsinesplineorder5}
\end{figure}


Figure \ref{figsinesplineorder5} shows the
effect of this modification on the estimation process for $N=6$, $\Delta t = .025$, 
$W=8$, $J=6$ and $K=5$.  The weight on the constraints was chosen as 
$ Q =$diag\{1.26, 5.19 10$^{-2}$, 1.54 10$^{-3}$, 4.4 10$^{-5}$, 1.3 10$^{-6}$,
4.4 10$^{-8}$\}.
The rationale for
selecting the entries of the
weight $Q$ is that the magnitude of each successive derivative grows by a
approximatley
30 (recall (\ref{scales}), or see (\ref{sines})), 
and thus $Q$ corresponds to equal ``relative
weighting'' on $y$ and its derivatives;
this same rationale will be applied to subsequent examples. Including the model
interpolation constraints allowed the same window size and delay as in Figure
\ref{figsineinterp} to be used with an increased polynomial order, without
``following'' the noise.


The method just outlined for estimating $x$ from $y$ and its derivatives is akin to
a ``dead-beat'' observer for $x$, though, as explained
earlier, there is smoothing in the
interpolation of $y$, which is ``IIR''.  Some additional smoothing from the
model can be obtained, and this constitutes the second method for introducing
the model into the interpolation/numerical differentiation process.
This technique, which is reminiscent of
\cite{grizzle1990b,moraal1992}, 
will be
directly applicable to nonlinear systems.  

Let $A_d$ denote the discretization of the linear model, (\ref{linearmodel}),
and let $\hat{x}_{k-1}^{+}(t_{k-1-W+K})$ be the previous estimate of $x$.  Update $x$
based on the latest estimate of $y$ and its derivatives through a damped Newton
method:
\begin{eqnarray}
\hat{x}_{k}^{-}(t_{k-W+K}) & = & A_d \hat{x}_{k-1}^{+}(t_{k-1-W+K}) \nonumber \\
\hat{x}_{k}^{+}(t_{k-W+K}) & = & \hat{x}_{k}^{-}(t_{k-W+K}) + 
      \epsilon P^{-1} \cdot \nonumber \\
   ( \left[ \begin{array}{l} 
\hat{y}_{k}(t_{k-W+K}) \\
\frac{d}{dt} \hat{y}_{k}(t_{k-W+K}) \\ 
\vdots \\ 
\frac{d^{n-1}}{dt^{n-1}} \hat{y}_{k}(t_{k-W+K}) \end{array} \right]& -& 
P  \hat{x}_{k}^{-}(t_{k-W+K}) ),\label{newton} 
\end{eqnarray}
where $P=col\{C, CA, \ldots, CA^{n-1})$ and $0 < \epsilon <  1$; using $\epsilon = 1$ reduces the above to a ``dead-beat''
estimator for $x$.  The damped Newton update can be used with
(\ref{leastsquares}) alone or with any of
the additions to (\ref{leastsquares}) discussed previously.

Selected simulations of the above estimation methods are now presented.
In what follows,
the state in (\ref{linearmodel}) was chosen as 
$$ x = col( y,  \frac{d}{dt} y ,  \frac{d^2}{dt^2} y , \frac{d^3}{dt^3} y).$$
It was properly initialized so that its output is precisely
given by (\ref{sines}).  The same noise sequence used in
Figure \ref{figsineinterp} was added to the output of (\ref{linearmodel}).
%also, the anti-aliasing filter (\ref{ysampled}) was used.


\begin{figure}[hbt]
\centerline{\psfig{file=sine_spline.eps,width=3.5in}}

\caption{\small Normalized estimation errors in \protect$y, \ \frac{dy}{dt}, \
\frac{d^2y}{dt^2}$} 
and  \protect$\frac{d^3y}{dt^3}$, respectively, incorporating 
``spline method'' and
Newton update in interpolation/numerical differentiation (equations
(\protect\ref{leastsquares}), (\protect\ref{derivconstraints}) and
(\protect\ref{modelconstraints})).
\label{figsinespline}
\end{figure}


Figure \ref{figsinespline} shows the normalized estimation error resulting
>from the application of (\ref{leastsquares})  in combination with
(\ref{derivconstraints})
and (\ref{newton}).   The parameter values were chosen as: $N=5$,
 $\Delta t = 0.025$ sec., $W=6$ and $K=3$  in (\ref{leastsquares}), 
$\eta = 4$ in (\ref{modelconstraints}) with $J=4$ and
a weight on (\ref{derivconstraints}) selected as
$Q =$diag\{1.26 10$^{-1}$, 5.19 10$^{-2}$, 1.54 10$^{-3}$, 4.4 10$^{-5}$\}
and $\epsilon = 0.03$ in (\ref{newton}).   
For comparison purposes,
a steady state, discrete-time, Kalman filter was designed with 
state noise covariance equal to 
$0.01 \ \sigma^2 \ Q^{-1}$ and output noise variance equal to $\sigma^2$; even
though (\ref{linearmodel}) does not have any process noise per se, the pair
consisting of the matrix $A$ and the state noise convariance matrix
had to be made stabilizable
for a stable Kalman filter to exist.
The magnitudes of the steady state errors were essentially identical to those in 
Figure \ref{figsinespline}.


%\begin{figure}[hbt]
%\centerline{\psfig{file=sine_model.eps,width=3.5in}}
%\caption{\small Normalized estimation errors in $y, \ \frac{dy}{dt}, \
%\frac{d^2y}{dt^2} \ $} 
%$\mbox{and} \ \frac{d^3y}{dt^3}$, respectively, incorporating model derivative
%constraints and damped Newton updates with 
%interpolation/numerical differentiation.
%\label{figsinemodel}
%\end{figure}


The normalized estimation errors resulting
>from the application of (\ref{leastsquares}) in combination with
(\ref{modelconstraints})
and (\ref{newton}) 
with  the parameter values chosen as: $N=5$, 
$\Delta t = 0.025$
sec., $W=6$ and $K=3$  in (\ref{leastsquares}), $\eta = 4$ in
(\ref{modelconstraints}), 
$J=5$ and a weight selected as
$ Q =$diag\{1.26, 5.19 10$^{-2}$, 1.54 10$^{-3}$, 4.4 10$^{-5}$, 1.3 10$^{-6}$\} were, again, essentially identical to those in 
Figure \ref{figsinespline}.


\begin{thebibliography}{11}

\small
\bibitem{diop1992b}
S.~Diop.
\newblock Rational system equivalence, and generalized realization theory.
\newblock In {\em Proceedings of the {I}FAC-{S}ymposium {N}OLCOS'92}, pages
  402--407, Bordeaux, 1992. IFAC.

\bibitem{diop1991d}
S.~Diop and M.~Fliess.
\newblock On nonlinear observability.
\newblock In C.~Commault, D.~Normand-Cyrot, J.~M. Dion, L.~Dugard, M.~Fliess,
  A.~Titli, G.~Cohen, A.~Benveniste, and I.~D. Landau, editors, {\em
  Proceedings of the {F}irst {E}uropean {C}ontrol {C}onference}, pages
  152--157, Paris, 1991. Herm{\`e}s.

\bibitem{UMTechReport93}
S. Diop, J.\ W.\ Grizzle, 
P.\ E.\ Moraal and A. Stefanopoulou, ``Interpolation and Numerical 
Differentiation for Observer Design'', University Michigan's College of
Engineering Control Group Report Series,  Report No. CGR-93-14, September
1993.


\bibitem{grizzle1990b}
J.~W. Grizzle and P.~E. Moraal.
\newblock Newton observers and nonlinear discrete-time control.
\newblock In {\em Proceedings of the 29th {I}EEE {C}onference on {D}ecision and
  {C}ontrol}, 1990, pages 760--767.

\bibitem{moraal1992}
P.~E. Moraal and J.~W. Grizzle.
\newblock Nonlinear discrete-time observers using {N}ewton's and {B}royden's
  methods.
\newblock In {\em Proceedings of the {A}merical {C}ontrol {C}onference}, June
1992, pages 3086--3090.

\bibitem{Sch73}
U. Sch\"{o}nwandt, ``Approximations to nonlinear observers'', {\it Automatica},
Vol. 9,1973,  pages 349--356.

\end{thebibliography}

\end{document}





